% !TEX root = ../main.tex
\chapter{Declaració d'ús d'intel·ligència artificial}\label{app:ai}
\fancyhead[RE]{\slshape \textsf{Apèndix \thechapter.\ Ús d'IA}}
\fancyhead[LO]{\slshape \textsf{\nouppercase \rightmark}}

\section{Objectiu de la declaració}

Aquest apèndix declara l'ús d'eines d'intel·ligència artificial generativa en
la preparació del projecte i de la memòria. La finalitat és deixar clar en
quines parts s'ha utilitzat suport automàtic, quins controls s'han aplicat i
quina responsabilitat manté l'autor sobre el resultat final.

L'eina principal utilitzada ha estat OpenAI Codex, emprada com a assistent de
programació, redacció i revisió. Aquest ús ha estat instrumental: l'eina ha
ajudat a inspeccionar el repositori, proposar fragments de codi o text,
organitzar comprovacions i detectar incoherències, però no ha substituït
l'execució dels experiments ni la interpretació final dels resultats.

Les xifres de la memòria no provenen d'una resposta generada per llenguatge.
Provenen de les execucions del pipeline C, de CSMR, de Raspberry i dels
extractes públics indicats a l'Apèndix~\ref{app:evidence}. Les comandes de
reproducció i verificació es documenten a l'Apèndix~\ref{app:reproducibility},
i la configuració de modes i presets a l'Apèndix~\ref{app:config}.

\section{Parts assistides amb IA}

La Taula~\ref{tab:ai_assisted_parts} resumeix les parts del treball on s'ha
utilitzat suport d'IA generativa. La columna final indica com s'ha controlat
cada ús abans d'incorporar-lo a la memòria o al repositori.

\begin{table}[H]
\centering
\small
\begin{tabularx}{\textwidth}{p{3.0cm}X X}
\toprule
Àmbit & Ús de la IA & Control aplicat \\
\midrule
Codi C i Python & Ajuda per revisar estructura, proposar scripts
d'orquestració i detectar incoherències entre rutes C, CSMR i Raspberry. &
Compilació, execució de tests, lectura manual del codi i comparació amb
artefactes generats. \\
Anàlisi de resultats & Suport per agregar CSV, revisar fórmules, connectar
taules amb checkpoints i preparar extractes públics. & Verificació amb
fitxers font, població, unitat, baseline i registre d'evidències. \\
Redacció \LaTeX & Propostes de redacció, reestructuració de capítols,
millora de captions i organització d'apèndixs. & Revisió humana, compilació
del PDF i decisió final de redacció segons el criteri de l'autor. \\
Figures i taules & Ajuda per decidir quines figures o taules explicaven millor
els resultats i per documentar-ne la lectura. & Les figures provenen de dades
generades o extractes públics; les captions s'han revisat manualment. \\
Revisió de coherència & Cerca de contradiccions sobre \code{q204},
\code{adaptive\_s8}, CSMR, Raspberry, Q-map i taula precalibrada. & Validació
creuada amb els capítols, apèndixs A--C i fitxers del repositori. \\
\bottomrule
\end{tabularx}
\caption{Parts del projecte assistides amb IA generativa i controls aplicats
abans d'incorporar-les a la versió final.}
\label{tab:ai_assisted_parts}
\end{table}

\section{Parts no substituïdes per IA}

L'ús de la IA no substitueix les parts essencials del treball. En particular:

\begin{itemize}[leftmargin=*]
  \item els experiments s'han executat amb el compressor, el descompressor, els
  generadors de Q-map, CSMR i els scripts de Raspberry;
  \item els valors de bitrate, PSNR, temps i cost computacional provenen de
  fitxers i extractes identificats, no de prediccions de l'assistent;
  \item les decisions metodològiques finals, l'abast del treball i la lectura
  crítica dels resultats corresponen a l'autor;
  \item les fonts bibliogràfiques s'han d'entendre com a responsabilitat de
  l'autor, encara que la IA hagi ajudat a revisar-ne la coherència editorial;
  \item la defensa del projecte, les limitacions i les conclusions són
  responsabilitat acadèmica de l'autor.
\end{itemize}

També cal distingir aquesta declaració de l'objecte tècnic del projecte.
SORTENY és un model neuronal de compressió d'imatges; no és una IA generativa de
llenguatge. Els presets automàtics implementats per a la selecció de regions són
regles deterministes sobre bandes i índexs espectrals, no decisions preses per
un assistent conversacional.

\section{Mesures de verificació}

Per reduir el risc d'errors, al·lucinacions o afirmacions no traçables s'han
aplicat les mesures següents:

\begin{itemize}[leftmargin=*]
  \item compilar la memòria fins eliminar referències, cites i figures sense
  resoldre;
  \item executar comprovacions locals com \code{git diff --check}, tests C i
  compilació del PDF;
  \item separar explícitament resultats mesurats, proxies C-only i treball
  futur;
  \item no acceptar una xifra sense unitat, població, baseline i font pública o
  origen arxivat;
  \item publicar extractes CSV petits per verificar les xifres principals sense
  necessitat de checkpoints complets;
  \item documentar a l'Apèndix~\ref{app:reproducibility} com executar proves
  bàsiques i com comprovar els extractes públics;
  \item mantenir a l'Apèndix~\ref{app:evidence} la traçabilitat entre
  afirmacions, fitxers públics i origen experimental.
\end{itemize}

Aquestes mesures no eliminen la necessitat de revisió humana. La IA pot
proposar una explicació convincent però incorrecta, confondre un proxy amb una
mètrica final o avançar conclusions abans que estiguin demostrades. Per això
les parts quantitatives de la memòria s'han vinculat a dades verificables i no
a text generat.

\section{Reflexió crítica}

L'ús de la IA ha tingut un impacte positiu en tres aspectes. Primer, ha accelerat
la navegació per un repositori gran, amb codi C, Python, figures, informes i
checkpoints. Segon, ha ajudat a convertir feedback tècnic extens en canvis
estructurats de la memòria. Tercer, ha facilitat revisions de coherència que són
costoses manualment, com mantenir estable el baseline \code{q204}, separar
\code{adaptive\_s8} com a control tècnic i no confondre CSMR amb el contenidor C.

El risc principal és que l'assistent tendeixi a omplir buits amb formulacions
plausibles. En aquest projecte això podia portar a tres errors especialment
greus: atribuir a la taula precalibrada una funció que no tenia, presentar
proxies C-only com a bitrate real, o afirmar viabilitat a bord més enllà de les
proves disponibles. Aquest risc s'ha mitigat exigint traçabilitat, compilació,
execució de scripts i revisió manual de les afirmacions.

També hi ha un impacte sobre l'autoria. La IA pot millorar la forma i suggerir
alternatives, però no pot assumir la responsabilitat del criteri tècnic. En
aquesta memòria, les contribucions del treball, la interpretació dels resultats
i la decisió de què és una evidència suficient corresponen a l'autor. La IA s'ha
utilitzat com a eina de suport, no com a substitut de l'aprenentatge, la
validació ni la responsabilitat acadèmica.
